\roadmapSection{4}{PROGRAMMING FOR HARDWARE AND SOFTWARE}{15--17 Weeks Total}

{\small\itshape\color{arligray} C is the foundation. Every other language builds on understanding C first. Learn C deeply, then branch.}

% ── 4.1 C Programming ─────────────────────────────────────────────────────────
\roadmapSubSection{4.1}{C Programming --- Foundation for Everything}{8 Weeks}

\roadmapHead{Language Fundamentals}
\checkitem{Variables and types: \texttt{char}, \texttt{int}, \texttt{unsigned int}, \texttt{long}, \texttt{float}, \texttt{double} --- understand sizes}{}
\checkitem{Understand that \texttt{sizeof(int)} varies by platform --- always check}{}
\checkitem{Operators: arithmetic, bitwise (\texttt{\& | \^{} \~{} << >>}), logical, comparison}{}
\checkitem{Control flow: if/else, for, while, do-while, switch --- and why \texttt{goto} exists in kernel code}{}
\checkitem{Functions: declaration, definition, calling convention, return values}{}
\checkitem{Understand stack vs heap: local variables on stack, \texttt{malloc} on heap}{}

\roadmapHead{Pointers --- Master Completely}
{\small\itshape\color{arligray} Maps directly to ARLIZ Vol I --- Chapter on Pointers. Write the same day you learn this.}

\checkitem{A pointer is an integer that holds a memory address. \texttt{int *p = \&x} means p holds address of x}{}
\checkitem{Dereferencing: \texttt{*p} reads or writes the value at the address stored in p}{}
\checkitem{Pointer arithmetic: \texttt{p + 1} advances by \texttt{sizeof(*p)} bytes, not 1 byte}{}
\checkitem{Pointer to pointer: \texttt{int **pp} --- used for output parameters and dynamic 2D arrays}{}
\checkitem{Function pointers: \texttt{void (*fn)(int)} --- used for callbacks and vtables}{}
\checkitem{void pointer: \texttt{void *} --- generic pointer, cast before use}{}
\checkitem{const correctness: \texttt{const int *p} vs \texttt{int * const p} --- know the difference}{}

\roadmapHead{Arrays and Memory}
{\small\itshape\color{arligray} Maps to ARLIZ Vol III --- Arrays chapter. Write while learning.}

\checkitem{Array name decays to pointer to first element in most contexts}{}
\checkitem{Array indexing: \texttt{a[i]} is exactly \texttt{*(a + i)} --- equivalent expressions}{}
\checkitem{Multidimensional arrays: row-major layout, how 2D maps to 1D memory}{}
\checkitem{String: array of \texttt{char} terminated by null byte --- understand \texttt{strlen} vs \texttt{sizeof}}{}
\checkitem{Stack allocation: \texttt{char buf[256]} is on stack --- know the risks of stack overflow}{}

\roadmapHead{Structs, Dynamic Memory, Build System}
\checkitem{Struct definition and member access via \texttt{.} and \texttt{->}}{}
\checkitem{Memory layout of struct: sequential, with alignment padding inserted by compiler}{}
\checkitem{\texttt{malloc(n)} allocates n bytes on heap, returns \texttt{void*} --- always check for NULL}{}
\checkitem{\texttt{free(p)} releases memory --- freeing NULL is safe, freeing twice is undefined behavior}{}
\checkitem{\texttt{calloc(n, size)} allocates and zeroes --- use for arrays}{}
\checkitem{Common bugs: use after free, double free, buffer overflow, memory leak}{}
\checkitem{Compilation stages: preprocessing $\to$ compilation $\to$ assembly $\to$ linking}{}
\checkitem{GCC basic usage: \texttt{gcc -o output source.c -Wall -Wextra}}{}
\checkitem{Write a basic Makefile: targets, prerequisites, recipes}{}
\newpage

\roadmapHead{Advanced C for Hardware Work}
\checkitem{\texttt{volatile} keyword: prevents compiler from caching hardware register reads --- required for MMIO}{}
\checkitem{Inline assembly \texttt{\_\_asm\_\_}: understand syntax, register constraints, clobber list}{}
\checkitem{Read an ELF binary: \texttt{.text} (code), \texttt{.data} (init globals), \texttt{.bss} (zero-init)}{}
\checkitem{Use \texttt{objdump -d} to disassemble compiled C and read the output}{}
\checkitem{Understand undefined behavior: signed overflow, null pointer deref, strict aliasing}{}
\checkitem{Implement ring buffer: head/tail pointer arithmetic, wrap-around, empty vs full}{}
\checkitem{Implement simple memory allocator: free list, split and coalesce blocks}{}

\newpage
% ── 4.2 Arduino and Microcontrollers ──────────────────────────────────────────
\roadmapSubSection{4.2}{Arduino and Microcontrollers (Talken Platform)}{3--4 Weeks}

\roadmapHead{Arduino Fundamentals}
\checkitem{Understand what Arduino is: AVR or ARM microcontroller + bootloader + USB interface + IDE}{}
\checkitem{Install Arduino IDE and understand its structure: \texttt{setup()} and \texttt{loop()}}{}
\checkitem{Digital pins: \texttt{INPUT}, \texttt{OUTPUT}, \texttt{INPUT\_PULLUP} modes}{}
\checkitem{\texttt{analogRead()}: 10-bit ADC, reads 0--5V as 0--1023}{}
\checkitem{\texttt{analogWrite()}: PWM output, 8-bit, $\approx$490Hz frequency}{}
\checkitem{\texttt{millis()} and \texttt{micros()} for non-blocking timing (never use \texttt{delay()} in real code)}{}
\checkitem{Serial: \texttt{begin(9600)}, \texttt{print()}, \texttt{println()}, \texttt{read()} --- UART under the hood}{}

\roadmapHead{Arduino Hardware Projects --- Build Each One}
\checkitem{Blink LED with resistor calculation: V=5V, LED Vf=2V, I=20mA $\to$ R = (5-2)/0.020 = 150$\Omega$}{}
\checkitem{Button input with debounce: read state, compare to previous, trigger on edge only}{}
\checkitem{PWM LED brightness control via potentiometer: \texttt{analogRead} $\to$ \texttt{map} $\to$ \texttt{analogWrite}}{}
\checkitem{Servo control: Servo library, understand 50Hz PWM and pulse width encoding}{}
\checkitem{Temperature sensor (DHT22 or DS18B20): read via library, understand protocol}{}
\checkitem{I2C LCD display: Wire library, understand I2C addressing}{}
\checkitem{Ultrasonic distance sensor (HC-SR04): measure pulse width for distance}{}
\checkitem{Serial communication between two Arduinos: TX of one to RX of other}{}

\roadmapHead{ESP32 Microcontroller --- Talken Platform}
\checkitem{ESP32 vs Arduino: dual-core ARM Cortex-M, built-in WiFi + BT, 240MHz}{}
\checkitem{Set up ESP-IDF development environment (professional toolchain, not Arduino IDE)}{}
\checkitem{ESP32 memory: IRAM, DRAM, external PSRAM, SPI Flash regions}{}
\checkitem{Write GPIO driver using ESP-IDF registers, not Arduino abstractions}{}
\checkitem{Write UART driver: configure clock divider, baud rate, TX/RX pins via GPIO matrix}{}
\checkitem{Write I2C driver: configure SCL/SDA, understand clock stretching}{}
\checkitem{Write SPI driver: configure MOSI/MISO/CLK/CS, understand DMA transfers}{}
\checkitem{Implement FreeRTOS tasks: create task, set priority, communicate via queue}{}
\checkitem{Implement ESP32 deep sleep: configure wakeup source, measure current in sleep mode}{}

\roadmapHead{Other Languages --- Learn After C Is Solid}
\checkitem{Python: \texttt{pyserial} for UART, \texttt{RPi.GPIO} for Raspberry Pi, \texttt{matplotlib} for data analysis}{}
\checkitem{Rust: ownership model eliminates use-after-free, \texttt{embedded-hal} traits, \texttt{no\_std}}{}
\checkitem{Bash: build automation, pipes/grep/awk/sed for log parsing, Makefile targets}{}
\checkitem{JavaScript/TypeScript: Node.js for serial port automation, ARLIZ documentation tools}{}

\newpage
% ── 4.3 Embedded C Patterns for Microcontrollers ─────────────────────────────
\roadmapSubSection{4.3}{Embedded C Patterns for Microcontrollers}{2--3 Weeks}

\roadmapHead{Memory-Mapped I/O (MMIO)}
\checkitem{Peripheral registers are memory addresses. Write to address = write to hardware register}{}
\checkitem{\texttt{volatile uint32\_t *GPIO\_OUT = (uint32\_t *)0x3FF44004;} then \texttt{*GPIO\_OUT |= (1<<pin)}}{}
\checkitem{Why volatile: compiler must not cache register reads. Every access goes to hardware}{}
\checkitem{Bit field structs: represent register layout as struct with bit fields. Read-modify-write}{}
\checkitem{Atomic operations: single instruction read-modify-write. \texttt{GPIO.out\_w1ts.val} on ESP32}{}

\roadmapHead{State Machine Patterns}
\checkitem{Enum state type: defines all valid states. Switch on current state}{}
\checkitem{Event-driven FSM: event struct passed to state handler, returns new state}{}
\checkitem{Hierarchical FSM: states contain sub-states. Exit/enter actions on transition}{}
\checkitem{Apply to: button debounce FSM, UART frame parser FSM, protocol state machine}{}
\checkitem{Build: UART command parser as FSM. States: IDLE, RECEIVING, PROCESSING, RESPONDING}{}

\roadmapHead{Bit Manipulation for Embedded}
\checkitem{Set bit n: \texttt{reg |= (1U << n)}}{}
\checkitem{Clear bit n: \texttt{reg \&= \~{}(1U << n)}}{}
\checkitem{Toggle bit n: \texttt{reg \^{}= (1U << n)}}{}
\checkitem{Test bit n: \texttt{(reg >> n) \& 1U}}{}
\checkitem{Extract bitfield: \texttt{(reg >> START) \& ((1U << WIDTH) - 1)}}{}
\checkitem{Insert bitfield: \texttt{reg = (reg \& mask) | (value << START)}}{}
\checkitem{Practice: decode a CAN frame, extract PGN from J1939 identifier using bitfield ops}{}

\roadmapHead{Circular Buffer Implementation}
\checkitem{Ring buffer: fixed-size array, head/tail indices, wrap at size using modulo}{}
\checkitem{Empty: head == tail. Full: (head + 1) \% size == tail. Never fill all slots (ambiguity)}{}
\checkitem{Thread-safe ring buffer: disable interrupts or use atomic tail/head updates}{}
\checkitem{DMA-backed ring buffer: DMA writes to buffer, CPU reads. Zero-copy UART receive}{}
\checkitem{Implement: 256-byte ring buffer in C. Test corner cases: empty read, full write, wrap}{}

\roadmapHead{Checksum and CRC}
\checkitem{Parity: count 1-bits, check even or odd. Detects 1-bit error only}{}
\checkitem{LRC (Longitudinal Redundancy Check): XOR all bytes. Simple but weak}{}
\checkitem{CRC-8: 8-bit polynomial checksum. Used in 1-Wire, SMBus}{}
\checkitem{CRC-16 CCITT: 16-bit. Used in XMODEM, X25, DNP3}{}
\checkitem{CRC-32: used in Ethernet, ZIP, PNG. IEEE 802.3 polynomial}{}
\checkitem{Implement CRC-16 in C: table lookup (fast) or bit-by-bit (small flash)}{}
\newpage
% ── 4.4 Python for Hardware Automation ───────────────────────────────────────
\roadmapSubSection{4.4}{Python for Hardware Automation and Testing}{2 Weeks}

\checkitem{pyserial: read and write UART. ser = serial.Serial('/dev/ttyUSB0', 115200, timeout=1)}{}
\checkitem{Parse binary protocol: struct.unpack('<HBB', data) for little-endian uint16 + 2 bytes}{}
\checkitem{Automated test: send command, read response, compare to expected. Log pass/fail}{}
\checkitem{matplotlib: plot ADC samples over time. Real-time plot with animation.FuncAnimation}{}
\checkitem{scipy.signal: apply Butterworth filter to recorded samples in Python}{}
\checkitem{numpy: array operations on measurement data. FFT with numpy.fft.fft()}{}
\checkitem{smbus2 (Linux): read I2C device directly from Raspberry Pi. No microcontroller needed}{}
\checkitem{RPi.GPIO: control GPIO from Raspberry Pi. Bit-bang protocols, sensor reading}{}
\checkitem{Build: Python script that reads ESP32 serial data, plots real-time sensor graph}{}